\chapter{The Human Body: A Brief Introduction}
\label{ch:human-body}

\emph{How is the body built and how do its systems work, so that a yoga teacher can explain what the practices act on?}
A yoga instructor is not a doctor, but must know the body well enough to explain a posture, a breath or a caution in plain terms. This chapter goes system by system: what each system is made of and what it does, with the facts a candidate is most often asked. Learn one function for each system, and learn the autonomic nervous system well, because it links stress and relaxation.

\section{The body as connected systems}

The body works as a set of connected systems, none of which works alone. The usual list is skeletal, muscular, nervous, endocrine, circulatory, respiratory, digestive, excretory and reproductive; some books add the immune or lymphatic system. The sense organs, the eyes, ears, nose, tongue and skin, bring information to the nervous system.

Together the systems keep the inner state of the body steady: its temperature, water and chemistry. This is \emph{homeostasis}. Yoga practice supports it through rest, slow breathing and a calm mind.

% ycb1-fig-bodysystems: the screenshot is retired; retyped here from the notes Material bullets.
\begin{table}[H]\centering\small
\begin{tabular}{@{}L{0.2\linewidth}L{0.36\linewidth}L{0.36\linewidth}@{}}
\toprule
\textbf{System} & \textbf{Main parts} & \textbf{Main work} \\
\midrule
Skeletal & Bones, cartilage, ligaments, tendons, joints & Framework, support and protection; movement with the muscles; blood cells formed in the marrow; mineral store \\
Muscular & Voluntary muscles of the body; involuntary muscles of organs such as the heart and gut & Moves the body and the organs \\
Nervous & Brain, spinal cord, nerves; the sense organs feed it & Control and quick messages; the autonomic part works without our will \\
Endocrine & Glands & Hormones in the blood that control growth, mood and energy \\
Circulatory & Heart, blood vessels & Carries blood, oxygen, nutrients and waste \\
Respiratory & Nose, windpipe, lungs; diaphragm and rib muscles & Takes in oxygen, gives out carbon dioxide \\
Digestive & The gut & Breaks down food, takes in nutrients, removes solid waste \\
Excretory & Kidneys and bladder; skin and lungs & Removes waste; balances water and salts \\
Reproductive & Reproductive organs & Continuation of the species \\
\bottomrule
\end{tabular}
\par\smallskip{\footnotesize\itshape The body systems at a glance.}
\end{table}

\section{The skeletal system}

\subsection*{Parts and functions}
The skeleton is made of bones, cartilage, ligaments, tendons and joints. Bones are living tissue with a good blood supply; an adult has about 206. Cartilage is elastic and cushions the joints; ligaments join bone to bone; tendons join muscle to bone; and joints are where bones meet. The skeleton gives the body its framework and support, protects the organs, makes movement possible together with the muscles, forms blood cells in the bone marrow, and stores minerals.

\subsection*{Divisions}
The \emph{axial} skeleton is the skull, the vertebral column and the rib cage, which guards the heart and lungs. The \emph{appendicular} skeleton is the shoulder girdle (shoulder blades and collar bones), the arm bones (humerus, radius, ulna, and the bones of the wrist, hand and fingers), the pelvis, and the leg bones (femur, patella, tibia, fibula, and the bones of the ankle, foot and toes).

\ycbfig{ycb1-fig-skeleton}
\ycbfig{ycb1-fig-skeleton-divisions}

\subsection*{The vertebral column}
The spine is a column of vertebrae joined by cartilage and ligaments, with the spinal cord running through its central canal. From the top, its regions are cervical (seven vertebrae), thoracic (twelve), lumbar (five), sacral (fused) and coccygeal (fused). Its natural curves are cervical, thoracic, lumbar and sacral (pelvic). Faulty posture can exaggerate them: \emph{kyphosis} is too much thoracic curve, \emph{lordosis} too much lumbar curve, and \emph{scoliosis} a sideways curve.

\ycbfig{ycb1-fig-vertebral-column}

\subsection*{Movements of the spine}
\begin{itemize}
\item \textbf{Flexion}, bending forward: paschimottanasana, padahastasana.
\item \textbf{Extension}, bending back: bhujangasana, dhanurasana.
\item \textbf{Rotation}, twisting: ardha matsyendrasana. Rotation is greatest between the first two neck vertebrae, the atlas and the axis.
\item \textbf{Side bending}: trikonasana.
\item \textbf{Elongation}, stretching upward: tadasana, urdhva hastasana.
\end{itemize}

\ycbfig{ycb1-fig-spinal-movements}

\subsection*{Joints}
Fibrous joints allow almost no movement, as in the sutures of the skull. Cartilaginous joints are strong and slightly flexible, as where the ribs meet the breastbone, which helps breathing. Synovial joints are the most mobile: their bone ends are coated with cartilage inside a capsule filled with fluid, and they come in several kinds for different movements.

\ycbfig{ycb1-fig-joints}

\section{The muscular system}

\subsection*{Types and functions}
There are three types of muscle. \emph{Voluntary} (skeletal) muscle is consciously controlled, for walking, balance and writing. \emph{Involuntary} (smooth) muscle works under autonomic control, in digestion and blood flow. \emph{Cardiac} muscle is found only in the heart and beats by its own rhythm. Muscles turn chemical energy into movement. They also hold posture against gravity, protect the organs, produce heat, store energy, and drive the working of the internal organs.

\ycbfig{ycb1-fig-muscle-chart-showing-the-muscular-system-labeled}

\subsection*{How a muscle contracts}
Each muscle fibre holds thin filaments (actin) and thick filaments (myosin), which slide over each other and shorten the fibre. In an \emph{isotonic} contraction the muscle changes length under a steady load; in an \emph{isometric} contraction it keeps its length while its tension rises, as in holding a posture.

Muscles burn fuel in two ways. \emph{Aerobic} work uses oxygen and suits steady, long activity such as dance or distance running. \emph{Anaerobic} work uses little oxygen and suits fast, heavy effort; it forms lactic acid. After hard effort the body takes in extra oxygen to restore its energy stores and clear the lactic acid; this extra oxygen is the oxygen debt.

\ycbfig{ycb1-fig-mechnaism-of-muscle-contrtaction}

\subsection*{Reflexes and reciprocal inhibition}
Proprioceptors in the muscles and joints report the position of the body to the brain and help posture and coordination. The stretch reflex is a quick contraction when a muscle is suddenly lengthened, which is one reason to move slowly into a stretch. In \emph{reciprocal inhibition}, when one muscle contracts its opposite relaxes, as the triceps relaxes when the biceps contracts; this gives smooth movement.

\ycbfig{ycb1-fig-reciprocal-inhibition}

\subsection*{Words for movement}
Flexion closes the angle of a joint and extension opens it. Abduction moves a limb away from the midline and adduction towards it. Elevation and depression move a part up and down, as in the shoulders. Pronation turns the palm down and supination turns it up. Rotation turns a part about its axis, and sphincter action narrows or widens an opening. The matching muscle groups are flexors, extensors, abductors, adductors and rotators (see the sukshma vyayama figures in \cref{ch:vyayama}).

\section{The nervous system}

The neuron is the basic unit of the nervous system. Its two main parts are central and peripheral.

\subsection*{Central nervous system}
The \emph{brain} has the cerebrum; the cerebellum, for balance and coordination; the medulla, for vital functions; and the diencephalon with the thalamus and hypothalamus. It is guarded by the skull and the meninges. The \emph{spinal cord} runs from the medulla to the lumbar region; it carries sensory and motor signals and controls reflexes.

\ycbfig{ycb1-fig-overview-nervious-system}
\ycbfig{ycb1-fig-basic-brain-anatomy}
\ycbfig{ycb1-fig-spinal-chord}

\subsection*{Peripheral nervous system}
The \emph{somatic} part is voluntary: sensory and motor nerves, in cranial and spinal pairs. The \emph{autonomic} part is involuntary and works without our will. Its \emph{sympathetic} branch prepares the body for action, ``fight or flight''; its \emph{parasympathetic} branch calms it, ``rest and digest''. This pair is the link between stress and relaxation, and the reason slow breathing and relaxation are used in stress management (see \cref{ch:health-promotion}).

\ycbfig{ycb1-fig-peripheral-nervous-system}
\ycbfig{ycb1-fig-somatic-nervous-system}
\ycbfig{ycb1-fig-autonomic-nervous-system}

\section{The endocrine system}

Hormones are chemical messengers released by glands into the blood. They act on target organs, are needed in small amounts, and disturb health when there is too much or too little. Some are water-soluble (proteins and amines) and some lipid-soluble (steroids). The guidebook names the pituitary, thyroid, parathyroid and pineal glands.

The hypothalamus controls the pituitary through releasing and inhibiting hormones. The pituitary, called the master gland, releases growth hormone, FSH and LH, prolactin (milk), TSH (for the thyroid), ACTH (for the adrenals), ADH (water retention) and oxytocin (contraction of the uterus). ADH and oxytocin are made in the hypothalamus and released from the posterior pituitary. The pineal gland makes melatonin, linked to the sleep rhythm; the thyroid makes thyroxine and calcitonin, for metabolism and calcium; and the parathyroids keep the calcium balance.

\ycbfig{ycb1-fig-endocrine-system}
\ycbfig{ycb1-fig-endocrine}
\ycbfig{ycb1-fig-pituitary-gland}

\section{The circulatory system}

\subsection*{Blood}
Plasma is the largest share of the blood. Red cells carry oxygen and carbon dioxide; they are made in the bone marrow and live about 120 days. White cells give immunity, and platelets stop bleeding by plugging damaged vessels.

\subsection*{Vessels and heart}
Arteries carry blood away from the heart; they branch into arterioles and then capillaries, where exchange takes place; blood returns through venules and veins. The \emph{systemic} circulation serves the body and the \emph{pulmonary} circulation serves the lungs. The heart has two atria, which receive blood, and two ventricles, which pump it out. Valves stop backflow: the tricuspid on the right and the bicuspid or mitral on the left. The system transports, protects through immunity, and keeps the inner balance.

\ycbfig{ycb1-fig-cardiovascularsystem}
\ycbfig{ycb1-fig-heart-diagram}

\section{The respiratory system}

\subsection*{The airway}
The nose filters, warms and moistens the air and gives the sense of smell. Air passes through the pharynx, the larynx (voice box), and the trachea (windpipe, held open by rings of cartilage), into the bronchi and bronchioles, and reaches the alveoli, where gases are exchanged. The right lung has three lobes and the left has two. A lining of mucus traps dust. The sinuses are air spaces near the nose that can become blocked.

\ycbfig{ycb1-fig-respiratory-system}

\subsection*{The muscles of breathing}
The diaphragm is a dome below the lungs that separates the chest from the abdomen. The intercostal muscles between the ribs lift the rib cage on the in-breath and lower it on the out-breath. Accessory muscles of the neck help in clavicular breathing, and the abdominal muscles drive a forced out-breath, as in kapalabhati. This is why the sectional breaths of \cref{ch:pranayama} work part by part. The respiratory system exchanges gases, keeps the acid-base balance of the blood, and makes speech possible.

\ycbfig{ycb1-fig-breathing}
\ycbfig{ycb1-fig-article-respiration-muscles-wc1}

\section{The digestive system}

The mouth chews the food, and saliva starts the digestion of carbohydrates. The oesophagus is only a passage. The stomach churns the food with acid and starts the digestion of protein; it absorbs almost nothing. The small intestine (duodenum, jejunum and ileum) does most of the digestion and absorption, and its villi widen the surface. The large intestine absorbs water and passes through the ascending, transverse and descending colon to the rectum and anus. Food moves along the tract by peristalsis, waves of muscular contraction.

\ycbfig{ycb1-fig-digestive-system}

\section{The excretory system}

The kidneys filter the blood through their nephrons; the ureters carry urine down by peristalsis; the bladder stores it; and the urethra lets it out. The skin and the lungs also remove waste. The system removes waste, regulates the volume and pressure of the blood, and controls its salts and acidity.

\ycbfig{ycb1-fig-human-excretory-system}

\section{The reproductive system}

The reproductive system is concerned with the continuation of the species. The male parts are the testes, which make sperm, the scrotum, the seminal vesicles, the prostate and the penis. The female parts are the ovaries, which make eggs, the oviducts or fallopian tubes, where fertilisation takes place, the uterus, which nurtures the embryo, and the vagina. Fertilisation gives a zygote, which grows into an embryo and then a foetus.

\ycbfig{ycb1-fig-reprosys}
\ycbfig{ycb1-fig-human-male-reproductive-system-en}
\ycbfig{ycb1-fig-female-reproductive-system}

\begin{recall}
\begin{itemize}
\item Systems: skeletal, muscular, nervous, endocrine, circulatory, respiratory, digestive, excretory, reproductive.
\item About 206 adult bones; spine: cervical, thoracic, lumbar, sacral, coccygeal.
\item Muscle: voluntary, involuntary, cardiac; reciprocal inhibition gives smooth movement.
\item Autonomic: sympathetic activates, parasympathetic calms.
\item Pituitary: master gland. Pineal: melatonin. Mitral valve: left side of the heart.
\item Breathing muscles: diaphragm, intercostals, neck accessories, abdominals.
\end{itemize}
\end{recall}

\subsection*{Worked examples}
\begin{itemize}
\item \textbf{Model answer: why does slow breathing calm a person?} Slow, deep breathing and relaxation favour the parasympathetic branch of the autonomic nervous system, which calms the body (``rest and digest''), against the sympathetic branch, which prepares it for action.
\item \textbf{Model answer: name the spinal movement in each posture.} Padahastasana, flexion; bhujangasana, extension; trikonasana, side bending; ardha matsyendrasana, rotation; tadasana, elongation.
\item \textbf{A student asks why a teacher says ``never bounce into a stretch''.} A sudden lengthening of a muscle triggers the stretch reflex, a quick contraction that works against the stretch; moving slowly avoids it.
\end{itemize}

\subsection*{Check yourself}
\begin{enumerate}
\item Name the body systems in the usual list.
\item What is homeostasis?
\item How many bones does an adult have, about?
\item Name the regions of the vertebral column.
\item What is reciprocal inhibition?
\item Which part of the nervous system calms the body?
\item Which gland is called the master gland?
\item Which hormone does the pineal gland make?
\item Which valve lies between the left atrium and the left ventricle?
\item How long does a red cell live?
\item Name the muscles of breathing.
\item Which spinal movement does bhujangasana show?
\end{enumerate}

\subsection*{Answers}
\begin{enumerate}
\item Skeletal, muscular, nervous, endocrine, circulatory, respiratory, digestive, excretory, reproductive.
\item Keeping the body's inner state (temperature, water, chemistry) steady.
\item About 206.
\item Cervical, thoracic, lumbar, sacral and coccygeal.
\item When one muscle contracts, its opposite relaxes, as with the biceps and triceps.
\item The parasympathetic branch of the autonomic system.
\item The pituitary.
\item Melatonin.
\item The bicuspid (mitral) valve.
\item About 120 days.
\item Diaphragm, intercostals, accessory neck muscles, and the abdominal muscles for forced exhalation.
\item Extension (back bend).
\end{enumerate}

\paragraph{Sources.} The official guidebook for Level I (2016, older scheme), chapters on anatomy and physiology and on asana; standard references for the bone count and vertebral regions; a published review on the life span of red blood cells (NIH, PMC8006275); OpenStax, \emph{Anatomy and Physiology} 2e (joints, muscle, the pituitary, the vertebral curves); a review of yoga and the autonomic nervous system (PMC3447533); the Board syllabus for Level 1.
